\chapter*{Abstract}
\addcontentsline{toc}{chapter}{Abstract}

This thesis presents a comprehensive, top-down engineering approach to the design and control of a modular robotic
limb for a hexapod platform. The primary objective is to overcome the performance limitations inherent in low-cost
commercial actuation, which typically hinder precise legged locomotion.

The research begins with the kinematic modeling and digital simulation of the limb in the MATLAB/Simulink environment,
establishing a simulation model to validate trajectory planning and inverse kinematics algorithms. The core contribution
of this work is the development of a novel, custom "Smart Actuator" module designed to replace standard hobbyist
servo-motors. This custom embedded solution integrates a high-resolution magnetic encoder, a high-current H-bridge
driver, and an STM32 microcontroller. Crucially, it implements a non-linear Switching PID control strategy specifically
tuned to mitigate stick-slip friction phenomena and mechanical backlash.

The system is validated through a Hardware-in-the-Loop (HIL) architecture, bridging the simulation environment with the
physical prototype. Experimental benchmarking demonstrates a dramatic improvement in performance.